\section{Label Names}\label{sec:Label}
In chapter \tqfullvref{sec:CompoundStatements} and again in chapter \tqfullvref{sec:LabelsAndBreakStatements} we discussed the use of labels\index{label}, but we said nothing about the label identifier.

Basically, a label\index{label} is tagging a code block, and therefore you should choose something as the label name that describes that code block.

As a label\index{label} is only visible inside the current scope, you need not to be very creative: you can use the label \lstinline|ForeverLoop| for each non-terminating loop in your code.

And of course it makes sense to use the suffix \bdq{Loop} for all loop labels\index{label}.

The only fixed rule is that you should use \textit{CamelCase}, with a capital letter as the first letter, for the labels\index{label}. 

\subsubsection{Principles}
\begin{enumerate}[label=P\arabic*.]
	\item{Use a label\index{label} to give a descriptive name to a code block.}

	\item{The label\index{label} name has to be CamelCase, beginning with a capital letter.
	
	Only letters are allowed, digits are not recommended.}
	
	\item{The suffix \bdq{Loop} for all loop labels\index{label} is recommended.}
\end{enumerate}

%==============================================================================
% End of file!!
%==============================================================================
